\section{Introduction}
\citet{milano2019moderate} found that environmental variation across generations improved evolved controllers, with the best performance at a moderate rate of change. Ecological neuroevolution imposes a different selection pressure: agents must acquire enough food to survive and reproduce. Food acquisition also determines task performance, but selection through survival may favor different controllers from direct optimization of food intake.

We compare ecological neuroevolution (OE), evolution strategies (ES) and proximal policy optimization (PPO) in a continuous two-dimensional foraging world. OE maintains an interacting population with energy-dependent death and reproduction, supplemented by a minimum population of externally supplied replacements. ES and PPO optimize food acquisition and deploy copies of a single controller. Each method has separately tuned fixed- and varied-environment configurations. We evaluate food acquisition before and during adaptation, then measure population persistence in a separate assay without replacements.

A supporting experiment with inherited traits and fixed behavioral rules examines search objectives, mutation and cloning (Appendix~\ref{app:trait}).

We predicted that varied-environment OE would outperform PPO with domain randomization during adaptation. OE instead starts below PPO-DR and remains below it throughout the measured interval. Varied OE also scores below fixed OE; a later ES comparison favors the varied configuration by a smaller amount.

\section{Related work}
\citet{milano2019moderate} varied the rate of environmental change during evolution. \citet{milano2024qualitative} compared ES and PPO across tasks and treatments of environmental variation. Our experiment compares separately selected fixed- and varied-training configurations in a foraging task without sweeping the rate of environmental change.

\citet{hamon2023eco} studied persistent neural-agent populations with physiology-dependent reproduction and death. \citet{nisioti2022plasticity} used a minimal ecological model and ablations to distinguish fitness-based selection from niche-limited competition. Our implementation also couples survival and reproduction to agent behavior, but uses feedforward controllers and externally supplied population replacements.

Prior comparisons examine curricula \citep{nisioti2025transfer}, sequential changes and forgetting \citep{nisioti2026continual}, and update and inheritance regimes in ecological evaluations \citep{ge2026arena}. Here, adaptation begins independently from the trained snapshot in each test environment and permits further parameter updates.

